Within an interval state, the earliest available rank must be assigned
to one candidate $p$. Choosing $p$ divides the remaining candidates
into left and right intervals, and the unused ranks must split between
those two subproblems. The recursion enumerates these partitions;
memoization reuses each interval-and-rank state.

\begin{algorithm}[H]
\caption{\SlotDP: exact interval deadlines with a rank subset}
\label{alg:slot-dp}
\KwIn{Ordinary interval demands and deadlines; region $[l,r]$;
available absolute ranks $T\subseteq[L]$.}
\KwOut{A feasible rank assignment, or infeasible.}
\If{no original demand is wholly contained in $[l,r]$}{\KwRet the empty assignment\;}
\If{$T=\varnothing$}{\KwRet infeasible\;}
\If{this state is memoized}{\KwRet its stored result\;}
$t\gets\min T$\;
\If{a wholly contained demand has deadline $<t$}{\KwRet infeasible\;}
\ForEach{$p\in[l,r]$ and $T_L\subseteq T\setminus\{t\}$}{
  $T_R\gets T\setminus(T_L\cup\{t\})$\;
  Recurse on $([l,p-1],T_L)$ and $([p+1,r],T_R)$\;
  \If{both assignments are feasible}{
    Memoize their union with $(p,t)$ and \KwRet it\;
  }
}
Memoize and return infeasible\;
\end{algorithm}

After the recursion succeeds, sort the selected candidates by rank and
suppress unused ranks to obtain the schedule. For circular demands,
enumerate the first candidate and cut the circle there; this converts
the remaining demands to intervals, with one subtracted from their
deadlines and the slot cap to account for the forced first rank.
In the optional complete-exit variant, first compute both oriented
decoder costs using Algorithm~\ref{alg:endpoint-code}. Those costs
replace the ordinary deadline precheck by the local height tests of
Section~\ref{sec:polynomial-slots}; the same rank partitions then
handle the remaining scheduling choices.
